%% This LaTeX-file was created by <root> Thu Jan  9 17:18:47 1997
%% LyX 0.10 (C) 1995 1996 by Matthias Ettrich and the LyX Team

%% Don't edit this file unless you are sure what you are doing.
\documentclass[10pt,a4paper,oneside,onecolumn]{article}
\usepackage[]{fontenc}
\usepackage[latin1]{inputenc}
\usepackage[dvips]{epsfig}

%%
%% BEGIN The lyx specific LaTeX commands.
%%

\makeatletter
\def\LyX{L\kern-.1667em\lower.25em\hbox{Y}\kern-.125emX\spacefactor1000}%
\newcommand{\lyxtitle}[1] {\thispagestyle{empty}
\global\@topnum\z@
\section*{\LARGE \centering \sffamily \bfseries \protect#1 }
}
\newcommand{\lyxline}[1]{
{#1 \vspace{1ex} \hrule width \columnwidth \vspace{1ex}}
}
\newenvironment{lyxbibliography}
{
\begin{thebibliography}{99}}
{\end{thebibliography}}
\def\lxq{"}
\newenvironment{lyxcode}
{\list{}{
\rightmargin\leftmargin
\raggedright
\itemsep 0pt
\parsep 0pt
\ttfamily
}%
\item[]
}
{\endlist}
\newcommand{\lyxlabel}[1]{#1 \hfill}
\newenvironment{lyxlist}[1]
{\begin{list}{}
{\settowidth{\labelwidth}{#1}
\setlength{\leftmargin}{\labelwidth}
\addtolength{\leftmargin}{\labelsep}
\renewcommand{\makelabel}{\lyxlabel}}}
{\end{list}}
\newcommand{\lyxletterstyle}{
\setlength\parskip{0.7em}
\setlength\parindent{0pt}
}
\newcommand{\lyxaddress}[1]{
\par {\raggedright #1 
\vspace{1.4em}
\noindent\par}
}
\newcommand{\lyxrightaddress}[1]{
\par {\raggedleft \begin{tabular}{l}\ignorespaces
#1
\end{tabular}
\vspace{1.4em}
\par}
}
\newcommand{\lyxformula}[1]{
\begin{eqnarray*}
#1
\end{eqnarray*}
}
\newcommand{\lyxnumberedformula}[1]{
\begin{eqnarray}
#1
\end{eqnarray}
}
\makeatother

%%
%% END The lyx specific LaTeX commands.
%%

\pagestyle{plain}
\setcounter{secnumdepth}{3}
\setcounter{tocdepth}{3}
\begin{document}


\subsection{Einleitung}

Ich bin mal wieder dabei, eine Idee für eine Quest zusammenzustellen.
Nachdem ich meine Ideen auch nur bedingt selber habe, dachte ich mir,
es wäre mal wieder an der Zeit, eine kleine Buchvorlage auszuweiten
(zumindest von der Idee her) und daraus eine kleine Quest zu basteln,
die ein wenig schwerer sein soll und auch für Buchkenner nicht einfach
zu lösen.

Dazu habe ich mir folgendes überlegt:


\section{Die Vier-Elemente-Quest}


\subsection{Die Questidee}

Wie der Name schon sagt, ist das Ding relativ einfach aufgebaut, auf
ein Thema, das in so vielen Fantasywerken schon behandelt und weitmöglichst
ausgebeutet wurde: die 4 Elemente, also Luft, Feuer, Erde, Wasser.
Darauf möchte ich meine Quest aufbauen, man kann sie für die, die
es kennen, vielleicht ein wenig mit Tarko Reub (aus dem MG von meinem
dortigen Papi Wargon) vergleichen, die ja auch darauf basiert, dass
man die 4 Elemente zu einem zentralen Punkt zusammenträgt und sich
daraufhin etwas tut.

Anhand des Vierteilers (wo leider der letzte Band noch nicht erhältlich
ist *hmpflgrrr*) hatte ich nun folgende Idee, Artefakte jedes Elements
zu schaffen, eine Harfe der Winde, ein Flammenschwert, einen Stab
der Erde und einen Kelch des Wassers. Jedes dieser Artefakte ist nun
natürlich nicht so leicht zu erringen, und schön langsam drängt sich
die Frage auf, ob man daraus nicht eine 4-Teile-Quest basteln könnte,
die aufeinander aufbaut. Zusätzlich hat nun jedes dieser Artefakte
eine ganz besondere Eigenschaft, die man (im aufbauenden Sinn) dazu
einsetzen muß, um die restlichen Dinge und questrelevanten Dinge zu
kriegen.


\subsection{Die Erschaffung der Artefakte}

Die alten Mächte, die lange vor den Wesen, die sich in unserer Welt
jetzt so herumtreiben, haben die Artefakte (der Macht, das ist jetzt
ein Originalzitat aus dem Buch) erschaffen, um die Elemente zu beherrschen
und um einen Weltuntergang, der allen prophezeit wurde, zu verhindern.

Dieser Weltuntergang sollte vor allem deswegen kommen, weil die meisten
Wesen mit magischen Fähigkeiten diese ja nicht nur zu ihrem besten
ausnutzten sondern wie überall gab es Neider und Geizhälse sowie die
Gier nach Macht.


\subsubsection{Die Harfe der Winde}

Die Harfe der Winde wurde dazu geschaffen, das Wetter zu beherrschen
und die Winde zu bändigen. Darunter darf man sich jetzt nicht vorstellen
(übrigens so wie bei allen andren Artefakten auch), daß man als Besitzer
dieses Dings sich einfach hinstellen kann und Gott spielen (wir haben
ja schon genug Götter ;)), sondern daß man dadurch bloß in den Lauf
der Dinge einen kleinen Einfluß nehmen kann, jedoch am Gesamtgefüge
nichts wesentliches modifiziert. Das Wetter geht weiterhin seinen
gewohnten Gang, ja, rächt sich sogar für unrechtmäßigen Gebrauch der
Harfe.

Nehmen wir doch als Beispiel, daß die Harfe während dem Winter eingesetzt
wird, um frühlingshafte Verhältnisse sprungartig herbeizuführen. Nach
den ersten Erfolgen werden Stürme und gewaltige Unwetter über das
Land ziehen, da das unmittelbare absichtliche Herbeiführen einer neuen
Jahreszeit natürlich alle natürlichen Vorgänge über den Haufen wirft
und diese erst mal wieder in die normalen Bahnen gelenkt werden sollten.


\subsubsection{Das Flammenschwert}

Dieses Artefakt ist das dritte, das man besitzen wird. Es befähigt
den Träger zu unglaublicher Macht im Schwertkampf und ist sofort tödlich,
sobald es ein andres Wesen trifft. Nachdem dieses Artefakt einzigartig
ist (so wie eigentlich alle anderen auch) sollten die Quests nicht
gleichzeitig gespielt werden können, da bin ich mir aber noch nicht
so ganz im klaren.

Die durchaus schlagkräftigen Argumente dieses Schwerts sollten allerdings
so in der Art wie es im MG bei der Schneewittchenquest gemacht ist,
abgeschwächt werden, denn eine Waffe die alles tötet, was einem über
den Weg läuft ist in einem Mud etwas unbrauchbar.

Dabei fällt mir gerade ein, warum ist es eigentlich nicht möglich,
mit einer Waffe mehr als einen Schadenstyp anzurichten? Ich meine,
ein Flammenschwert ist doch erst mal a) ein Schwert, b) greift es
noch mit Feuer an, da es magisch ist auch noch mit c) DT\_MAGIC...
das ganze könnte man mit kombinierten Waffen noch ausweiten.


\subsubsection{Der Erdenstab}

Der Stab der Erde hat ganz einen simplen Zweck: die Beherrschung aller
Flora und Fauna, die auf Erden herumkreucht und -fleucht. Mit ihm
gewinnt man Macht über das Wachstum der Pflanzen, kann mit den Tieren
reden und ihnen befehlen und so weiter.

Leider fehlt mir auch hierfür eine passende Verwendung, aber ich wollte
schon weiter oben erwähnen, daß ich mich über jede Mitarbeit an diesem
schon etwas größeren Projekt (so wie es jetzt ausschaut) freue und
gerne auch andere Gegenden miteinbinden mag.


\subsubsection{Der Kessel des Wassers}

Das letzte der Artefakte und damit der Schlußpunkt dieser einen (vier)
Quest(s) birgt große Geheimnisse in sich, die Macht über das Leben.
Man kann ihn als hl. Gral auffassen, aber so drastisch würde ich in
der effektiven Realisierung im Silberland (oder wie auch immer) nicht
arbeiten, unsterblich sind wir eh schon alle (wir Magier sowieso,
die Spieler auch mit Einschränkungen in Form einer Property, die langsam
nach oben wandert...). Aber dafür wird es auch noch genug Ideen geben.


\subsection{Die Geschichte}

Zwar wuseln mir Ideen durch den Kopf, jedoch bin ich gern für Anregungen
jedweder Art offen.

(to be continued)


\subsection{Programmiertechnische Details}


\subsubsection{Räume}

Für die Räume denk ich da wieder mal an einen Prototyp pro Questteil,
der als eigene Klasse oder ähnlich implementiert wird, von dem dann
die einzelnen Räume abgeleitet werden.

Den Stab der Erde wird man wohl in einer Stadt namens Dhiammara finden,
diese befindet sich zwar in einem Berg, allerdings nicht so wirklich.
Näheres siehe Geschichte.

Die Harfe der Winde bekommt man bei ihrer Hüterin namens Cailleac'h,
diese lebt ebenfalls unterirdisch in einem steinernen Wald.

Das Flammenschwert kommt aus einem feuerroten Rubin der auf einer
Insel in einem Tal (mit See natürlich) befindet, nur der Questende
kann diesen Rubin sehen, alle anderen nicht.

Der Kessel der Harfe, als Krönung der Quest(s), tja, da hab ich mir
noch nix überlegt. Vielleicht als verstecktes Item in einer uralten
Bibliothek, in der das gesammelte Wissen von Jahrhunderten steckt.
Aber die Bücher mag ich nicht beschreiben :)).


\subsubsection{Objekte}

Ausser den Quest-Items und den Artefakten soll es natürlich den üblichen
Anteil an Waffen, Rüstungen und sonstigen Gegenständen geben, die
man halt so brauchen kann. Keine konkreten Vorstellungen.


\subsubsection{NPCs}

Tja, außer Info-NPCs und Quest-NPCs mit Aktionen und den üblichen
Metzel-NPCs gibts wohl eh kaum was anderes. Keine konkreten Vorstellungen


\subsubsection{Sonstiges}

Lästiger Punkt, aber den hebe ich mir auf. Das File wird schließlich
geupdatet.


\subsection{Randbemerkungen}

Natürlich sollten alle diese Gebiete als solche nicht brachliegen
und nur dem Quest-Zweck dienen. Dazu müssen die Gebiete alle gut gemacht
sein und die Zusatz-Features der entsprechenden Zonen Anreiz genug,
nicht nur die Quest zu machen. Allerdings gebe ich zu bedenken, daß
die Artefakte möglichst unique sein sollten... das legt wiederum nahe,
die Questgebiete vom übrigen Spielerbetrieb zu sperren. Eine blöde
Diskrepanz aber mit den nötigen Anregungen von euch bestimmt lösbar!

So, das solls fürs erste gewesen sein! Ich bitte um rege Mitarbeit
:)

\begin{quotation}

-Crea- / Michi

\end{quotation}

\end{document}
